% !TEX root = ../main.tex
% =====================================================================
% CAPITOLO 5 — VALUTAZIONE SPERIMENTALE
% (ex capitolo 7; in espansione su richiesta dei relatori — vedi
%  plans/PIANO_REVISIONE_FEEDBACK.md, attività F2)
% =====================================================================

\chapter{Valutazione sperimentale}
\label{chap:orion-valutazione}

Questo capitolo presenta la valutazione sperimentale del sistema su quattro piani complementari.
Il primo, e più esteso, è il terzo contributo della tesi: una metodologia di valutazione del routing di Orion tramite uno strumento di esecuzione automatica integrato nell'admin panel,
progettata per misurare in modo sistematico e ripetibile l'accuratezza delle decisioni dell'orchestratore e per guidare il miglioramento del \emph{system prompt} del router.
L'oggetto di questa iterazione --- il prompt stesso, con le tecniche di prompt engineering che incorpora --- è esaminato nella Sez.~\ref{sec:eval-prompteng};
i risultati su due versioni successive del prompt (v1.0 e v1.1) ne illustrano sia l'utilità sia i limiti.

Il secondo piano (Sez.~\ref{sec:eval-agenti}) valuta la \emph{qualità dei sub-agenti}: execution accuracy per gli agenti Text-to-SQL e faithfulness per il RAG,
applicando le metodologie introdotte nel background (Sez.~\ref{subsec:bg-text2sql} e~\ref{sec:bg-ragas}).

Il terzo piano (Sez.~\ref{sec:eval-costo-qualita}) analizza il trade-off costo--qualità realizzato dal routing per complessità,
chiudendo il percorso aperto nel background (Sez.~\ref{sec:bg-routing}); il quarto (Sez.~\ref{sec:eval-costi}) propone una valutazione economica indicativa
del costo di implementazione del progetto rispetto all'adozione di licenze SaaS.
Chiude il capitolo la discussione complessiva dei risultati e dei limiti dell'impianto sperimentale (Sez.~\ref{sec:eval-discussione}).

% =====================================================================
\section{Domande di ricerca}
\label{sec:eval-research-questions}
% =====================================================================

La valutazione è guidata da sei domande di ricerca:

\begin{description}
  \item[RQ1] Quanto è accurato il routing dell'orchestratore? (\emph{Agent Accuracy})
  \item[RQ2] Quanto è accurata la stima di complessità? (\emph{Complexity Accuracy})
  \item[RQ3] Il confidence score è calibrato? È più alto sui casi corretti che su quelli errati?
  \item[RQ4] Il meccanismo di clarification riduce l'errore senza degradare l'esperienza utente?
  \item[RQ5] Iterazioni successive del system prompt migliorano le metriche in modo misurabile?
  \item[RQ6] Il routing per complessità preserva la qualità delle risposte riducendo il costo? (\emph{costo--qualità})
\end{description}

La metodologia adottata combina elementi dell'approccio \emph{AgentBench}~\cite{liu2023agentbench}
per la valutazione di agenti LLM con il paradigma del prompt engineering iterativo~\cite{sahoo2024prompteng_survey}:
un gold standard fisso, un'esecuzione automatica, metriche dichiarate a priori, iterazione controllata del prompt con confronto quantitativo.

% =====================================================================
\section{Metodologia di valutazione}
\label{sec:eval-metodologia}
% =====================================================================

\subsection{Costruzione del gold standard}
\label{subsec:eval-gold-standard}

La test suite è composta da 120 test case, ciascuno definito da un prompt in linguaggio naturale, 
un agente atteso (\texttt{self}, \texttt{sysaid}, \texttt{allycare}, \texttt{coopolicy}), una complessità attesa (\texttt{low} o \texttt{high}) e una categoria tematica.

La costruzione è avvenuta in tre fasi:

\begin{enumerate}
  \item \textbf{Campionamento da log reali.} I prompt più frequenti nelle conversazioni di ciascun agente (raccolti durante la fase di test della suite) 
  sono stati anonimizzati e riformulati per eliminare riferimenti personali. 
  Questa fase ha prodotto circa 80~prompt rappresentativi del traffico effettivo.

  \item \textbf{Prompt sintetici per casi limite.} Sono stati creati manualmente circa 30~prompt progettati per esplorare le zone di ambiguità: 
  sovrapposizioni tra SysAid e AllyCare, tra Orion e CooPolicy (domande che potrebbero essere sia generiche sia documentali), 
  e richieste deliberatamente fuori dominio.

  \item \textbf{Prompt di sicurezza.} Un sottoinsieme di 10~prompt testa il comportamento del router su richieste malevole (prompt injection, richieste di dati personali) 
  per verificare che vengano gestite dalla pipeline di sicurezza.
\end{enumerate}

Le categorie tematiche sono cinque: \emph{SysAid} (36~casi), \emph{NPS} (21~casi), \emph{Policy} (33~casi), \emph{General} (23~casi) e \emph{Ambiguous} (7~casi).
La distribuzione riflette approssimativamente il volume di traffico atteso, con un \emph{oversampling} deliberato della categoria \emph{Ambiguous} per testare il meccanismo di clarification.

I test case sono raccolti in un file CSV (\texttt{test\_suite\_orion.csv}) e importati nel database dell'admin panel attraverso la pagina \emph{Eval Routing} (Sez.~\ref{subsec:eval-admin-page}).

\begin{figure}[htbp]
    \centering
    \includegraphics[width=\textwidth]{eval1}
    \caption[Estratto della test suite]{La test suite di valutazione nel tab ``Test Suite'' dell'admin panel di Orion: i 120 prompt etichettati con agente atteso, 
    complessità attesa e categoria tematica, importabili da CSV.}
    \label{fig:test_suite}
\end{figure}

\subsection{Architettura del framework di valutazione}
\label{subsec:eval-admin-page}

La valutazione avviene interamente dalla pagina \emph{Eval Routing} dell'admin panel di Orion (Sez.~\ref{sec:orion-impl-admin}), senza necessità di script CLI esterni o accesso diretto all'infrastruttura di Cooperativa.

\paragraph{Backend (\texttt{eval.py}).} Cinque endpoint REST esposti sotto il prefisso \texttt{/api/\allowbreak orion-admin/\allowbreak eval/}:

\begin{itemize}
  \item \texttt{GET /test-suite}: restituisce i test case attivi dal database.
  \item \texttt{POST /test-suite/import}: importa i test case da un contenuto CSV, sovrascrivendo i precedenti.
  \item \texttt{POST /run}: esegue la valutazione. Per ciascun test case, il backend chiama la funzione \texttt{classify()} \emph{direttamente} (invocazione Python, nessuna chiamata HTTP), 
  confronta l'output con l'atteso, salva il risultato nel database e costruisce il summary.
  \item \texttt{GET /runs}: restituisce la lista delle run passate con metriche aggregate.
  \item \texttt{GET /runs/\{run\_id\}}: restituisce il dettaglio completo di una run specifica.
\end{itemize}

L'invocazione diretta della funzione \texttt{classify()} è molto conveniente: 
il backend admin (\texttt{orion-admin-backend}) include una copia del modulo \texttt{router\_agent.py}, 
rendendo possibile l'import Python diretto senza la necessità di chiamare l'endpoint \texttt{/classify} del backend utente. 

Ogni test case costa una sola chiamata a Gemini Flash ($\sim$0,5~secondi); una run completa di 120~test case richiede circa 60~secondi.

\paragraph{Frontend (\texttt{Eval.tsx}).} Una pagina React organizzata in tre tab (\emph{Test Suite}, \emph{Lancia Eval}, \emph{Storico Run}) 
che integra l'intero ciclo di valutazione in un'unica interfaccia. Il tab \emph{Lancia Eval} visualizza i risultati con KPI card codificate cromaticamente 
(verde se il target è raggiunto, rosso altrimenti), breakdown per agente con barre orizzontali, 
lista dei casi errati espandibile con il reasoning del router, e tabella completa dei risultati con icone di correttezza.

\subsection{Metriche e target}
\label{subsec:eval-metriche}


\begin{table}[htbp]
  \centering
  \caption[Metriche di valutazione e target]{Metriche di valutazione del routing con target dichiarati a priori.}
  \label{tab:eval-metriche}
  \begin{tabular}{llc}
    \toprule
    \textbf{Metrica}                  & \textbf{Definizione}                              & \textbf{Target} \\
    \midrule
    Agent Accuracy                    & $\frac{\text{agente corretto}}{\text{totale}}$     & $\geq 85\%$ \\[3pt]
    Complexity Accuracy               & $\frac{\text{complessità corretta}}{\text{totale}}$ & $\geq 80\%$ \\[3pt]
    Clarification Rate                & $\frac{\text{clarification poste}}{\text{totale}}$  & $< 15\%$ \\[3pt]
    Mean Confidence (correct)         & Media della confidence sui casi corretti            & $\geq 0{,}82$ \\[3pt]
    Mean Confidence (wrong)           & Media della confidence sui casi errati              & $\leq 0{,}65$ \\
    \bottomrule
  \end{tabular}
\end{table}

Le ultime due metriche, prese congiuntamente, valutano la \emph{calibrazione} del confidence score~\cite{kadavath2022knowsknows}: 
un router ben calibrato è sicuro quando ha ragione e dubbioso quando sbaglia. 
Se le due distribuzioni di confidence (corretti vs errati) sono ben separate, la soglia di clarification può essere posizionata in modo efficace.

% =====================================================================
\section{Il prompt del router: anatomia e tecniche}
\label{sec:eval-prompteng}
% =====================================================================

Prima di presentare i risultati è utile esaminare l'oggetto su cui agisce
l'iterazione misurata in questo capitolo: il \emph{system prompt} del router.
Questa sezione ne descrive l'anatomia attraverso estratti della versione
v1.0, riprodotti dal codice pubblicato
(Sez.~\ref{sec:concl-generalizzabilita}, modulo
\texttt{orion-backend/app/agents/router\_agent.py}), e cataloga le tecniche
di prompt engineering che il prompt incorpora, collegandole alla letteratura
introdotta nel background (Sez.~\ref{sec:bg-prompting}); la stessa
catalogazione è poi estesa ai prompt degli agenti verticali, che
costituiscono altrettanti casi applicativi su dati aziendali reali.

\subsection{Anatomia del system prompt}
\label{subsec:eval-prompt-anatomia}

Il prompt è organizzato in blocchi delimitati da intestazioni maiuscole
--- una pratica di \emph{role prompting} con delimitatori espliciti che la
letteratura indica come efficace per ridurre le ambiguità di
parsing~\cite{sahoo2024prompteng_survey}. Il Listato~\ref{lst:prompt-regole}
riporta l'apertura e le regole di classificazione.

\begin{lstlisting}[caption={Estratto del system prompt del router (v1.0): ruolo, blocco agenti e regole di classificazione. Il segno [...] indica omissioni; i segnaposto fra graffe sono compilati a runtime.},label={lst:prompt-regole}]
Sei il modulo di routing di ORION, l'orchestratore AI
dell'organizzazione. Il tuo compito è classificare la richiesta
dell'utente per instradarlo all'agente specializzato più appropriato.

AGENTI DISPONIBILI:
{agents_block}

REGOLE DI CLASSIFICAZIONE:
1. Scegli l'agente il cui dominio corrisponde meglio alla richiesta
   dell'utente.
2. Se la richiesta è ambigua o potrebbe essere gestita da più agenti,
   scegli quello più probabile e indica una confidence più bassa.
3. Se la richiesta è chiaramente fuori dal dominio di tutti gli agenti
   specializzati, usa "self".
4. Saluti, conversazione generica, domande su documenti caricati
   dall'utente, e ricerche web vanno sempre a "self".
5. Per TUTTE le richieste, valuta la complessità in modo binario:
   - "low": domande dirette, conteggi semplici, filtri base, ricerche
     per un singolo criterio, aggregazioni semplici
   - "high": analisi multi-dimensionali, correlazioni, confronti
     temporali articolati, root cause analysis, report executive,
     ragionamento su più metriche incrociate
   La complessità influenza il modello LLM scelto dall'agente
   (es. SysAid usa un modello più potente per richieste complesse).
[...]
\end{lstlisting}

Le regole combinano tre dispositivi distinti. Le regole~1--3 definiscono la
politica di selezione e il comportamento di \emph{default} (la regola~4
assegna esplicitamente a \texttt{self} le classi di richieste più frequenti
e meno informative, sottraendole all'inferenza). La regola~5 definisce le
classi di complessità \emph{per esempi di istanze} anziché per definizione
astratta: è la stessa logica del few-shot
prompting~\cite{brown2020gpt3}, applicata però alla descrizione delle
etichette. La chiusura della regola~5 esplicita al modello la
\emph{conseguenza} della sua decisione (la selezione del modello nel
sub-agente): rendere espliciti gli effetti delle scelte è una pratica
raccomandata per i compiti di classificazione
strumentale~\cite{sahoo2024prompteng_survey}.

Il secondo estratto (Listato~\ref{lst:prompt-confidence}) riguarda la stima
della confidence, il fondamento del meccanismo di clarification
(Sez.~\ref{subsec:orion-confidence}).

\begin{lstlisting}[caption={Estratto del system prompt del router (v1.0): rubrica di confidence e istruzione di deferral.},label={lst:prompt-confidence}]
STIMA DELLA CONFIDENCE:
- 0.90-1.00: la richiesta è chiaramente nel dominio di un singolo
  agente, senza ambiguità
- 0.70-0.89: la richiesta è probabilmente nel dominio indicato, ma ci
  sono elementi di ambiguità
- 0.50-0.69: la richiesta è ambigua, potrebbe appartenere a più domini
- sotto 0.50: la richiesta è troppo vaga per essere classificata con
  sicurezza

Se la confidence è inferiore a {confidence_threshold}, imposta
clarification_needed con una domanda breve e chiara per disambiguare
l'intento dell'utente.
\end{lstlisting}

La confidence non è chiesta come numero libero ma ancorata a una
\emph{rubrica} di quattro fasce con descrizioni operative: è una forma di
\emph{verbalized confidence elicitation}, che la letteratura sulla
calibrazione indica come ragionevolmente affidabile nei modelli di grande
scala~\cite{kadavath2022knowsknows}. L'istruzione finale collega la stima
all'azione: sotto soglia il modello deve formulare esso stesso la domanda di
disambiguazione, realizzando il \emph{deferral} human-in-the-loop ispirato
all'active learning~\cite{settles2012active,diao2023activeprompting}.
Il campo \texttt{reasoning} obbligatorio dell'output (Chain-of-Thought
esplicita~\cite{wei2022cot}) e il vincolo di \emph{guided
generation} sullo schema JSON con \texttt{enum} dinamica
(Sez.~\ref{subsec:orion-output-strutturato},~\cite{willard2023outlines})
completano il quadro: il primo fornisce tracciabilità diagnostica, il
secondo elimina per costruzione i target inesistenti.

Un'ultima proprietà è operativa: il testo del prompt è un \emph{template}
con segnaposto (\texttt{\{agents\_block\}}, \texttt{\{history\_block\}},
\texttt{\{confidence\_threshold\}}, \texttt{\{user\_prompt\}}) compilati a
ogni richiesta, ed è persistito nel database (\texttt{app\_settings}),
modificabile dalla pagina \emph{Router Prompt} dell'admin panel
(Sez.~\ref{sec:orion-impl-admin}) senza rilascio di codice. È questa
proprietà a rendere praticabile il ciclo di iterazione misurato nella
Sez.~\ref{sec:eval-risultati}.

\subsection{Few-shot dinamico: il blocco degli agenti}
\label{subsec:eval-prompt-fewshot}

Il segnaposto \texttt{\{agents\_block\}} è compilato a runtime con la
descrizione dei soli agenti disponibili per l'utente corrente: per ciascuno,
nome, descrizione del dominio ed esempi di richieste tipiche. Il
Listato~\ref{lst:prompt-agente} riporta il blocco di SysAid.

\begin{lstlisting}[caption={Estratto del system prompt del router (v1.0): il blocco descrittivo di SysAid generato a runtime, con gli esempi few-shot.},label={lst:prompt-agente}]
### SysAid (Ticket IT) (key: "sysaid")
Agente specializzato nell'analisi dei ticket IT SysAid. Interroga
BigQuery per cercare, filtrare, aggregare e analizzare i ticket di
assistenza tecnica (service requests). [...] Il modello LLM viene
selezionato automaticamente in base alla complessità della richiesta.
Esempi di richieste tipiche:
  - "Quanti ticket sono stati aperti questa settimana?"
  - "Mostrami i ticket aperti per il negozio di Bologna"
  - "Analizza il trend dei ticket per categoria negli ultimi 6 mesi
     e identifica le anomalie"
  - "Genera un report executive sui KPI del service desk dell'ultimo
     trimestre"
\end{lstlisting}

Gli esempi sono \emph{in-context examples} nel senso di
Brown et al.~\cite{brown2020gpt3}, con due accortezze. La prima è che gli
esempi di ciascun agente sono scelti per coprire entrambe le fasce di
complessità: i primi due del listato modellano richieste \texttt{low}
(conteggio, filtro singolo), gli ultimi due richieste \texttt{high}
(trend con anomalie, report executive); il blocco insegna così
simultaneamente il dominio dell'agente \emph{e} la nozione di complessità
della regola~5. La seconda è che il blocco è \emph{dinamico}: riflette i
soli agenti abilitati per l'utente, in coerenza con l'\texttt{enum}
dell'output strutturato, così che descrizioni e vincoli di schema non
possano mai divergere.

\subsection{Le tecniche nei prompt della suite}
\label{subsec:eval-prompt-tecniche}

Il router non è l'unico componente in cui il comportamento del modello è
governato dal prompt: ciascun agente verticale incorpora tecniche analoghe,
descritte nel Capitolo~\ref{chap:agenti} insieme alla loro funzione. La
Tab.~\ref{tab:eval-tecniche} le riassume in chiave comparata, collegandole
alla letteratura di riferimento.

\begin{table}[htbp]
  \centering
  \caption[Tecniche di prompt engineering nella suite]{Le tecniche di prompt engineering utilizzate nella suite, con riferimento alla letteratura e punto di applicazione.}
  \label{tab:eval-tecniche}
  \small
  \begin{tabular}{p{3.4cm}p{2.6cm}p{7.2cm}}
    \toprule
    \textbf{Tecnica} & \textbf{Riferimento} & \textbf{Applicazione nella suite} \\
    \midrule
    Role prompting e sezioni delimitate & \cite{sahoo2024prompteng_survey}
      & Incipit e struttura a blocchi del router (Listato~\ref{lst:prompt-regole}); prompt di sistema ``collegiali'' degli agenti (Sez.~\ref{sec:agenti-principi}) \\
    In-context examples (few-shot) & \cite{brown2020gpt3}
      & Esempi di richieste tipiche per agente nel blocco agenti del router (Listato~\ref{lst:prompt-agente}) \\
    Chain-of-Thought esplicita & \cite{wei2022cot}
      & Campo \texttt{reasoning} obbligatorio del router; disciplina \texttt{\#\#\# RAGIONAMENTO}/\texttt{\#\#\# RISPOSTA} degli agenti (Sez.~\ref{subsec:agenti-auth}) \\
    Guided generation (output strutturato) & \cite{willard2023outlines}
      & \texttt{response\_schema} JSON con \texttt{enum} dinamica nel router (Sez.~\ref{subsec:orion-output-strutturato}); classificatore di corpus di CooPolicy (Sez.~\ref{subsec:coopolicy-multicorpus}) \\
    Confidence verbalizzata con rubrica & \cite{kadavath2022knowsknows}
      & Le quattro fasce della stima di confidence (Listato~\ref{lst:prompt-confidence}) \\
    Deferral / clarification & \cite{settles2012active,diao2023activeprompting}
      & \texttt{clarification\_needed} sotto soglia di confidence (Sez.~\ref{subsec:orion-confidence}) \\
    Schema injection selettiva e glossario semantico & \cite{pourreza2023dinsql}
      & AllyCare e SysAid: schema delle sole tabelle rilevanti, glossario delle colonne in italiano, anagrafica negozi iniettata (Sez.~\ref{subsec:allycare-schema}) \\
    Regole negative e di dominio & \cite{sahoo2024prompteng_survey}
      & Filtri di validità SQL e \texttt{UPPER(TRIM(...))} (Sez.~\ref{subsec:allycare-schema}); ``gerarchia della verità'' e fiducia nel tool (Sez.~\ref{subsec:allycare-discipline}); le otto regole non negoziabili di CooPolicy (Sez.~\ref{subsec:coopolicy-generation}) \\
    Self-correction & \cite{pourreza2023dinsql}
      & Errori SQL reimmessi come \texttt{function\_response} nel ciclo agentico (Sez.~\ref{subsec:allycare-tools}) \\
    \bottomrule
  \end{tabular}
\end{table}

Per diverse di queste tecniche l'effetto è riscontrabile negli esiti
documentati nel resto della tesi. La ``fiducia nel tool'' di AllyCare ha
eliminato in fase di test la classe di errori in cui il modello, ricalcolando
mentalmente, contraddiceva il risultato della propria query
(Sez.~\ref{subsec:allycare-discipline}). Le regole SQL di dominio si
riflettono nell'alta quota di SQL eseguibile prodotto dagli agenti
Text-to-SQL ($23/25$ e $25/25$ nelle run della Sez.~\ref{subsec:eval-text2sql}).
La guided generation con \texttt{enum} elimina per costruzione gli
instradamenti verso agenti inesistenti, riducendo il fallback difensivo a
salvaguardia residuale (Sez.~\ref{subsec:orion-impl-router}). La rubrica di
confidence, infine, è il presupposto della calibrazione verificata in
Sez.~\ref{subsec:eval-confidence-cal}.

Il prompt v1.0 qui descritto non è un primo tentativo: è l'esito delle
iterazioni informali condotte durante lo sviluppo della suite, guidate
dall'osservazione qualitativa dei log di routing.
% TODO F3.5: arricchire con 2-3 passate di tuning pre-v1.0 documentate
% (recuperare versioni storiche del prompt da app_settings / history
% Bitbucket e note di sviluppo).
L'iterazione successiva (v1.0~$\rightarrow$~v1.1), condotta invece con il
metodo quantitativo di questo capitolo, è documentata nella
Sez.~\ref{subsec:eval-iterazione}.

% =====================================================================
\section{Risultati}
\label{sec:eval-risultati}
% =====================================================================

\subsection{Run v1.0: primo prompt}
\label{subsec:eval-v10}

La prima run di valutazione (prompt version v1.0) ha prodotto i seguenti risultati su 120~test case (Tab.~\ref{tab:eval-v10}):

\begin{table}[htbp]
  \centering
  \caption[Risultati della run v1.0]{Risultati della run v1.0 (120 test case, primo system prompt).}
  \label{tab:eval-v10}
  \begin{tabular}{lcc}
    \toprule
    \textbf{Metrica}             & \textbf{Valore}  & \textbf{Target} \\
    \midrule
    Agent Accuracy               & $99{,}2\%$ (119/120) & $\geq 85\%$ \\
    Complexity Accuracy          & $88{,}3\%$ (106/120) & $\geq 80\%$ \\
    \bottomrule
  \end{tabular}
\end{table}

L'accuracy globale supera il target del $85\%$ fin dalla prima iterazione. Di seguito il breakdown per agente (Tab.~\ref{tab:eval-v10-breakdown}):

\begin{table}[htbp]
  \centering
  \caption[Breakdown per agente --- v1.0]{Agent accuracy per agente target nella run v1.0.}
  \label{tab:eval-v10-breakdown}
  \begin{tabular}{lcl}
    \toprule
    \textbf{Agente}   & \textbf{Accuracy}     & \textbf{Note} \\
    \midrule
    SysAid            & 100\% (36/36)         & Perfetto \\
    CooPolicy         & 100\% (33/33)         & Perfetto \\
    AllyCare          & 95\% (20/21)          & Eccellente \\
    Self (Orion)      & 100\% (30/30)         & Perfetto \\
    \bottomrule
  \end{tabular}
\end{table}

\begin{figure}[htbp]
    \centering
    \includegraphics[width=\textwidth]{eval2}
    \caption[Risultati della run v1.0 nell'admin panel]{La pagina di valutazione al termine della run v1.0: le quattro KPI card 
    (agent accuracy, complexity accuracy, confidence media sui casi corretti e sugli errati) e il breakdown per agente.}
    \label{fig:accuracy_per_agent}
\end{figure}

\subsection{Analisi degli errori della v1.0}
\label{subsec:eval-errori-v10}

L'errore di classificazione della v1 ricade in questo pattern: 

\paragraph{Pattern 1: Correlazione cross-dataset non riconosciuta (1~errore).} Il prompt ``Analizza la correlazione tra NPS e volume di vendita per punto vendita'' 
viene instradato a \texttt{self} anziché ad \texttt{allycare}, perché il router non riconosce che i dati NPS sono necessari per l'analisi. 
È l'unico errore, con confidence bassa (0,55), segnale che il router stesso percepiva l'incertezza.

\subsection{Iterazione del prompt: v1.1}
\label{subsec:eval-iterazione}

Sulla base dell'analisi sono state apportate queste modifiche, per cercare di eliminare l'errore e migliorare la complexity accuracy:

\begin{enumerate}
  \item \textbf{Regola 8: Correlazioni cross-dataset.} Aggiunta la regola che richieste che \emph{richiedono dati} di un agente specializzato (anche se formulate in modo generico) 
  vanno instradate a quell'agente.

  \item \textbf{Complessità.} È stata modificata la classificazione per cercare di fare correttamente catalogare come ``low'' le ricerche a singola query o i confronti semplici.

\end{enumerate}

Entrambe le modifiche sono istanze delle tecniche catalogate nella
Sez.~\ref{sec:eval-prompteng}: la Regola~8 è una \emph{regola di precedenza}
esplicita, che risolve i conflitti fra domini sovrapposti a favore
dell'agente che detiene i dati necessari; la revisione dei criteri di
complessità arricchisce la definizione per istanze delle classi
\texttt{low}/\texttt{high} (Listato~\ref{lst:prompt-regole}, regola~5) con i
casi di confine emersi dall'analisi degli errori.
% TODO F3.5: inserire qui il diff testuale effettivo v1.0 -> v1.1 del prompt
% (recuperare la v1.1 dalla tabella app_settings via pagina Router Prompt).

La run v1.1 è stata eseguita sullo \emph{stesso} gold standard della v1.0 modificando esclusivamente il system prompt del router. La Tab.~\ref{tab:eval-v11} confronta le due run.

\begin{table}[htbp]
  \centering
  \caption[Confronto fra le run v1.0 e v1.1]{Confronto fra le run v1.0 e v1.1 sul medesimo gold standard di 120 test case. Cambia esclusivamente il system prompt del router.}
  \label{tab:eval-v11}
  \begin{tabular}{lcc}
    \toprule
    \textbf{Metrica}                 & \textbf{v1.0} & \textbf{v1.1} \\
    \midrule
    Agent Accuracy                   & $99{,}2\%$      & $99{,}2\%$ \\
    Complexity Accuracy              & $88{,}3\%$      & $93{,}3\%$ \\
    Errori (su 120)                  & 1               & 1 \\
    Confidence media (corretti)      & $0{,}922$       & $0{,}922$ \\
    Confidence media (errati)        & $0{,}55$        & $0{,}5$ \\
    \bottomrule
  \end{tabular}
\end{table}

Il guadagno più netto è sulla \emph{complexity accuracy} ($+5{,}0$~punti percentuali): le regole aggiunte hanno reso più stabile la distinzione low/high. 
L'\emph{agent accuracy} invece non migliora: la Regola~8 non è bastata a correggere la correlazione cross-dataset dell'errore, che ricompare a bassa confidence. 


\begin{figure}[htbp]
  \centering
  \includegraphics[width=\textwidth,keepaspectratio]{eval4}
  \caption[Risultati della run v1.1 nell'admin panel]{Risultati della run v1.1 nell'admin panel, con lo storico delle due run sul medesimo gold standard. 
  Le KPI card e il breakdown per agente si riferiscono alla v1.1.}
  \label{fig:eval_v11}
\end{figure}

\subsection{Calibrazione del confidence score}
\label{subsec:eval-confidence-cal}

La calibrazione del confidence score è verificata confrontando la distribuzione della confidence sui casi corretti e quella sui casi errati. 
Un router ben calibrato presenta due distribuzioni ben separate: confidence alta sui casi corretti, confidence bassa sui casi errati.

Nella run v1.0 la confidence media è risultata pari a $0{,}922$ sui casi corretti e $0{,}55$ sul caso errato;
la seconda run l'ha ulteriormente abbassata a $0{,}5$.

% =====================================================================
\section{Valutazione della qualità degli agenti verticali}
\label{sec:eval-agenti}
% =====================================================================

La valutazione del routing (Sez.~\ref{sec:eval-risultati}) misura se Orion sceglie l'agente \emph{giusto}, ma non se l'agente scelto risponde \emph{bene}. 
Questa sezione colma quel divario valutando direttamente la qualità dei sub-agenti, con due metodologie distinte a seconda della forma di grounding: 
\emph{execution accuracy} per gli agenti Text-to-SQL (SysAid, AllyCare) e \emph{faithfulness} con \emph{answer relevancy} per l'agente RAG (CooPolicy). 
Entrambe sono applicazioni dirette delle metriche introdotte nel Capitolo~\ref{chap:background} (Sez.~\ref{subsec:bg-text2sql} e~\ref{sec:bg-ragas}).

Il codice della valutazione è pubblicato nel modulo \texttt{evaluation/} del repository (Sez.~\ref{sec:concl-generalizzabilita}): 
gli script invocano gli \emph{endpoint reali} degli agenti così che la misura sia end-to-end e non su una riproduzione parziale del sistema.

\subsection{Text-to-SQL: execution accuracy}
\label{subsec:eval-text2sql}

La metrica adottata è l'\emph{execution accuracy} (o \emph{denotation accuracy}), lo standard dei benchmark Spider~\cite{yu2018spider} e 
BIRD~\cite{li2023bird}: una query generata è corretta se, eseguita, produce lo \emph{stesso result-set} della query di riferimento (\emph{gold}), 
indipendentemente dalla sua formulazione testuale. Rispetto al confronto sintattico fra stringhe SQL, l'execution accuracy è la misura corretta perché query 
diverse possono essere semanticamente equivalenti.

\paragraph{Protocollo.} Per ciascuna domanda di un gold standard etichettato (domanda in linguaggio naturale + query SQL di riferimento) lo script invoca l'agente sul suo 
endpoint di produzione, estrae le query SQL effettivamente eseguite (esposte dall'agente nel campo \texttt{queries}),
le esegue su BigQuery insieme alla gold e ne confronta i result-set in modo order-insensitive e robusto a casing e arrotondamenti. 
Si riportano due metriche: \textbf{EX (final)}, che richiede il match con l'ultima query del turno (quella che produce la risposta finale), ed \textbf{EX (any)}, 
che ammette il match con almeno una delle query eseguite; la distanza fra le due separa gli errori di generazione SQL dagli errori di selezione della query conclusiva.

\begin{table}[htbp]
  \centering
  \caption[Execution accuracy degli agenti Text-to-SQL]{Execution accuracy degli agenti Text-to-SQL su un gold standard di domande etichettate con SQL di riferimento. La complessità segue la stessa distinzione low/high del routing.}
  \label{tab:eval-text2sql}
  \begin{tabular}{lcccc}
    \toprule
    \textbf{Agente} & \textbf{N} & \textbf{EX (final)} & \textbf{EX (any)} & \textbf{IC 95\% (Wilson)} \\
    \midrule
    AllyCare (NPS)      & 25 & 80{,}0\% & 80{,}0\% & $[60{,}9\% ,\, 91{,}1\%]$ \\
    SysAid (ticketing)  & 25 & 84{,}0\% & 84{,}0\% & $[65{,}3\% ,\, 93{,}6\%]$ \\
    \bottomrule
  \end{tabular}
\end{table}

AllyCare raggiunge un'execution accuracy dell'$80\%$ (20/25) e SysAid dell'$84\%$ (21/25), con SQL eseguibile prodotto rispettivamente in $23/25$ e $25/25$ dei casi. 
La coincidenza fra EX (final) ed EX (any) indica che, quando l'agente sbaglia, sbaglia la query stessa e non la mera selezione di quella conclusiva. 
Gli errori si concentrano su due classi note nella letteratura Text-to-SQL: (i)~le \emph{aggregazioni con regole di dominio implicite} come 
il calcolo dell'NPS (formula promotori/detrattori), le medie con filtri di validità (\texttt{voto\_item >= 0}) e le percentuali per AllyCare; 
le aggregazioni temporali per SysAid, dove l'interpretazione dell'agente diverge da quella codificata nella query gold; 
(ii)~alcuni casi in cui l'agente risponde in modo discorsivo senza eseguire SQL. Si osserva inoltre un effetto del \emph{non-determinismo}: 
un conteggio elementare classificato correttamente in un'esecuzione può fallire in un'altra, segnale che parte della varianza è stocastica e non sistematica 
(cfr. Sez.~\ref{subsec:eval-limiti}). Va infine ricordato che una quota di questi disaccordi è imputabile all'etichettatura del gold standard (un solo annotatore): 
su domande con regole di business sfumate la query di riferimento è una scelta progettuale. Gli intervalli di confidenza al $95\%$ (metodo di Wilson) sono ampi, $[60{,}9\%,\,91{,}1\%]$ per AllyCare
e $[65{,}3\%,\,93{,}6\%]$ per SysAid, conseguenza diretta del campione ridotto ($n=25$): i valori vanno letti come stime puntuali con incertezza non trascurabile.
Il breakdown di queste metriche per fascia di complessità, rilevante per il trade-off costo--qualità, è analizzato nella Sez.~\ref{subsec:eval-cq-breakdown}.

\subsection{RAG: faithfulness e answer relevancy}
\label{subsec:eval-rag}

Per CooPolicy si adottano due metriche del framework RAGAS~\cite{es2023ragas}, calcolate con il paradigma \emph{LLM-as-judge}~\cite{zheng2023llmjudge}: 
la \textbf{faithfulness} (frazione delle affermazioni della risposta supportate dai frammenti recuperati, misura anti-allucinazione) e l'\textbf{answer relevancy} 
(pertinenza della risposta alla domanda). Il gold standard distingue domande \emph{nel} corpus, su cui si misurano le due metriche e l'\emph{answer rate}, 
da domande \emph{fuori} corpus, su cui si misura il \emph{tasso di corretto rifiuto}: per un agente vincolato al solo materiale recuperato (Sez.~\ref{subsec:coopolicy-generation}), 
rifiutare correttamente una domanda fuori dominio è esso stesso un comportamento desiderato.

\begin{table}[htbp]
  \centering
  \caption[Metriche RAG di CooPolicy]{Metriche di qualità del RAG di CooPolicy, calcolate con LLM-as-judge (Gemini~2.5 Pro) su un gold standard di domande in-corpus e fuori-corpus.}
  \label{tab:eval-rag}
  \begin{tabular}{lc}
    \toprule
    \textbf{Metrica} & \textbf{Valore} \\
    \midrule
    Faithfulness media (in-corpus)              & 0{,}73 \\
    Answer relevancy media (in-corpus)          & 0{,}85 \\
    Answer rate (in-corpus)                     & 100\% (20/20) \\
    Corretto rifiuto (fuori-corpus)             & 80\% (4/5) \\
    Citazioni sulle risposte in-corpus          & 100\% (20/20) \\
    \bottomrule
  \end{tabular}
\end{table}

CooPolicy risponde a tutte le domande in corpus ($20/20$), sempre corredando la risposta di citazioni alle fonti, 
e rifiuta correttamente $4$ delle $5$ domande fuori corpus: l'unico scivolone è una richiesta di traduzione, 
a cui l'agente risponde dalla conoscenza parametrica anziché astenersi, comportamento innocuo ma formalmente fuori mandato. 
La faithfulness media ($0{,}73$) e l'answer relevancy ($0{,}85$) indicano risposte pertinenti e in larga parte ancorate al materiale recuperato, 
ma con \emph{varianza elevata}: diverse risposte sono pienamente fedeli (faithfulness $\geq 0{,}95$), 
altre contengono affermazioni non supportate dagli \emph{snippet} recuperati ($< 0{,}4$). 
Il valore va letto come stima \emph{conservativa}: il giudizio di grounding è calcolato sugli \emph{snippet} di 300~caratteri esposti dall'agente, 
non sul documento integrale, e penalizza quindi affermazioni corrette ma supportate da porzioni di documento non incluse nello snippet.

\paragraph{Limiti.} Le due valutazioni ereditano i limiti generali dell'impianto sperimentale (Sez.~\ref{subsec:eval-limiti}: dimensione del gold standard, bias del costruttore) e ne aggiungono di specifici.
L'execution accuracy dipende dalla correttezza delle query gold, etichettate da un singolo annotatore; il giudizio di faithfulness è affidato a un LLM, 
con i noti bias di posizione, lunghezza e auto-preferenza~\cite{zheng2023llmjudge}, e valuta il \emph{grounding} sugli snippet recuperati (300~caratteri) 
anziché sul documento integrale. Restano dunque misure indicative, utili a stabilire un \emph{ordine di grandezza} della qualità degli agenti e una base di partenza riproducibile, 
non un giudizio definitivo.

% =====================================================================
\section{Ottimizzazione costo--qualità}
\label{sec:eval-costo-qualita}
% =====================================================================

Il routing per complessità è il dispositivo con cui la suite realizza il
trade-off costo--qualità introdotto nel background
(Sez.~\ref{sec:bg-routing}), e opera su due livelli: all'interno degli
agenti, dove l'etichetta \texttt{low}/\texttt{high} seleziona il modello
(Gemini Flash o Pro per Vera, CooPolicy e AllyCare; Gemini Flash o Claude
Sonnet per SysAid, Sez.~\ref{subsec:vera-routing}); e nell'orchestratore,
dove il router stima la complessità e la propaga ai sub-agenti
(Sez.~\ref{subsec:orion-output-strutturato}). Rispetto ai sistemi di
riferimento, il meccanismo è una stima \emph{ex ante} a passaggio singolo:
non un cascade a tentativi successivi come
FrugalGPT~\cite{chen2023frugalgpt}, né un classificatore appreso su dati di
preferenza come RouteLLM~\cite{ong2024routellm}. Questa sezione ne esamina
i due lati: la struttura del risparmio (RQ6, lato costo) e il breakdown
della qualità per fascia di complessità (RQ6, lato qualità).

\subsection{Distribuzione della complessità e struttura del costo}
\label{subsec:eval-cq-costo}

Sul gold standard di 120 casi --- costruito per riflettere la composizione
del traffico atteso (Sez.~\ref{subsec:eval-gold-standard}) --- l'etichetta
di complessità è \texttt{low} in 104 casi ($86{,}7\%$) e \texttt{high} nei
restanti 16 ($13{,}3\%$). Indicando con $f_{\mathrm{low}}$ e
$f_{\mathrm{high}}$ le due frazioni e con $c_{\mathrm{low}}$ e
$c_{\mathrm{high}}$ i costi medi di una richiesta servita rispettivamente
dal modello economico e dal modello premium, il costo medio per richiesta
con routing è
\[
  c_{\mathrm{routing}} = f_{\mathrm{low}} \, c_{\mathrm{low}} + f_{\mathrm{high}} \, c_{\mathrm{high}},
\]
e il risparmio rispetto al servire ogni richiesta con il modello premium
vale
\[
  S \;=\; \frac{c_{\mathrm{high}}}{c_{\mathrm{routing}}}
    \;=\; \frac{\rho}{f_{\mathrm{low}} + f_{\mathrm{high}}\,\rho},
  \qquad \rho = \frac{c_{\mathrm{high}}}{c_{\mathrm{low}}},
\]
dove $\rho$ è il rapporto di costo fra i due modelli. Con la distribuzione
osservata, $S$ cresce con $\rho$ fino al limite asintotico
$1/f_{\mathrm{high}} = 7{,}5$: a titolo illustrativo, per $\rho = 5$ il
routing riduce il costo dell'inferenza di un fattore $3{,}3$, per
$\rho = 10$ di un fattore $4{,}5$, per $\rho = 20$ di un fattore $5{,}7$.
Il valore effettivo di $\rho$ dipende dai listini per token dei modelli e
dal volume medio di token per richiesta, misurabile dalla tabella
\texttt{chat\_logs} che registra i token in/out di ogni interazione
(Sez.~\ref{subsec:agenti-security}); la conversione del modello in valori
monetari è discussa fra i limiti (Sez.~\ref{subsec:eval-limiti}).
% TODO F3.3: token medi in/out per richiesta, per agente e modello
% (query su chat_logs in ambiente GCP).
% TODO F3.6: listini per Mtoken verificati (giugno 2026) di Gemini 2.5
% Flash/Pro su Vertex AI e Claude Sonnet 4.6 — calcolare rho effettivo e
% istanziare S. NON inserire numeri non verificati.
Al costo della risposta si somma quello del routing step di Orion, una
chiamata Flash per ogni richiesta, trascurabile per costruzione rispetto
alla generazione della risposta completa
(Sez.~\ref{subsec:orion-model-routing}).

% TODO F3.1: subsection "Il costo del router: Flash vs Pro" — run della
% test suite v1.1 con router su Gemini 2.5 Pro (etichetta "v1.1-pro")
% dall'admin panel: confrontare agent/complexity accuracy e costo per
% chiamata, a supporto sperimentale della scelta motivata in
% Sez. subsec:orion-model-routing e Tab. tab:orion-decisioni.

\subsection{Qualità per fascia di complessità}
\label{subsec:eval-cq-breakdown}

Il presupposto del routing è che il risparmio non degradi la qualità: le
richieste semplici devono essere servite adeguatamente dal modello
economico e quelle complesse dal modello premium. Una verifica si ottiene
disaggregando l'execution accuracy degli agenti Text-to-SQL
(Sez.~\ref{subsec:eval-text2sql}) per fascia di complessità del gold
standard --- la stessa etichetta che, attraverso la stima del
classificatore, determina il modello utilizzato. La
Tab.~\ref{tab:eval-cq-breakdown} riporta il breakdown, calcolato sulle
medesime run della Sez.~\ref{subsec:eval-text2sql}.

\begin{table}[htbp]
  \centering
  \caption[Execution accuracy per fascia di complessità]{Breakdown dell'execution accuracy per fascia di complessità del gold standard, con il modello che serve ciascuna fascia. Stesse run della Tab.~\ref{tab:eval-text2sql}.}
  \label{tab:eval-cq-breakdown}
  \small
  \begin{tabular}{llccc}
    \toprule
    \textbf{Agente} & \textbf{Fascia (modello)} & \textbf{N} & \textbf{EX (final)} & \textbf{IC 95\% (Wilson)} \\
    \midrule
    AllyCare & low (Gemini 2.5 Flash) & 17 & $82{,}4\%$ (14/17) & $[59{,}0\%,\,93{,}8\%]$ \\
    AllyCare & high (Gemini 2.5 Pro)  & 8  & $75{,}0\%$ (6/8)   & $[40{,}9\%,\,92{,}9\%]$ \\
    SysAid   & low (Gemini 2.5 Flash) & 16 & $75{,}0\%$ (12/16) & $[50{,}5\%,\,89{,}8\%]$ \\
    SysAid   & high (Claude Sonnet 4.6) & 9 & $100\%$ (9/9)     & $[70{,}1\%,\,100\%]$ \\
    \bottomrule
  \end{tabular}
\end{table}

Due osservazioni. La prima: sulle richieste complesse la qualità non
degrada --- l'esito peggiore per un sistema con routing sarebbe
un'accuratezza bassa proprio sulla fascia \texttt{high}, e il campione
mostra l'opposto, con il code path Claude Sonnet di SysAid che non commette
errori ($9/9$). La seconda: gli errori si concentrano sulla fascia
\texttt{low} servita dal modello economico ($12/16$ per SysAid), dove si
era già osservata una componente di varianza stocastica
(Sez.~\ref{subsec:eval-text2sql}): alcune richieste semplici falliscono per
instabilità più che per incapacità sistematica. Il margine di miglioramento
più ampio, dunque, non sta nell'estendere l'uso del modello premium ma nello
stabilizzare il comportamento del modello economico sulle richieste di
routine, oppure nel rivedere il confine fra le due classi di complessità.

\paragraph{Limiti del breakdown.} Tre cautele nella lettura.
(i)~I campioni per fascia sono piccoli (8--17 casi) e gli intervalli di
Wilson conseguentemente ampi. (ii)~Le due fasce contengono domande diverse
per costruzione: il confronto non misura la differenza fra modelli a parità
di domanda (servirebbe un disegno controllato con le stesse domande su
entrambi i code path); misura però la qualità del sistema così come operato
dal routing, che è la grandezza rilevante per l'utente finale.
(iii)~L'attribuzione del code path avviene tramite l'etichetta del gold
standard: in modalità standalone l'agente stima localmente la complessità
(Sez.~\ref{subsec:allycare-tools}) e su una minoranza di casi la stima può
divergere dall'etichetta; il breakdown va quindi letto come un \emph{proxy}
del code path effettivamente percorso.

La prospettiva di questa sezione è il costo marginale per richiesta; la
sezione successiva completa il quadro economico spostandosi sul costo di
adozione della piattaforma.

% =====================================================================
\section{Valutazione economica}
\label{sec:eval-costi}
% =====================================================================

La motivazione di costo (Cap.~\ref{chap:introduzione}) merita una verifica quantitativa, sia pure indicativa:
si costruisce un modello \emph{parametrico} con assunzioni esplicite e facilmente aggiornabili;
il calcolatore è incluso nel modulo \texttt{evaluation/cost\_model.py}.
L'obiettivo non è una stima contabile precisa, ma il confronto fra due modelli di costo: quello di licenze SaaS nominali per l'intera popolazione
e quello di una piattaforma a consumo come quella realizzata.

\paragraph{Assunzioni.} Si considera la popolazione di dipendenti \emph{potenzialmente abilitabili} a una licenza: in Coop Alleanza si tratta di circa $N \approx 4\,000$ figure con mansioni d'ufficio. 
Le suite SaaS di AI generativa per l'enterprise
(assistenti integrati nella suite di produttività o chatbot dedicati) hanno, a listino pubblico nel 2026, un prezzo dell'ordine di \texteuro20--30 per utente al mese
\footnote{Il riferimento è ai prezzi di listino pubblici dei principali assistenti enterprise (es. \texteuro30/utente/mese per le suite di produttività integrate), 
convertiti per omogeneità con il resto della tesi. Il valore è un parametro del modello, non un dato dell'organizzazione.}; nel modello si adotta il valore mediano di \texteuro25.
La piattaforma realizzata, al contrario, ha un costo dominato dal consumo: chiamate ai modelli su Vertex AI,
esecuzione serverless su Cloud Run (che scala a zero) e servizi managed di base. 
La differenza strutturale è che il SaaS è tariffato \emph{per utente abilitato}, la piattaforma \emph{per utilizzo effettivo}, ed è proprio su questo divario
che si fonda il vantaggio economico.

Il punto cruciale è che \emph{avere accesso} a uno strumento di AI non equivale a \emph{usarlo}. I dati di adozione lo confermano: secondo Gallup~\cite{gallup2026aiwork},
fra i lavoratori in ruoli d'ufficio (\emph{remote-capable}) solo il $40\%$ usa l'AI almeno qualche volta a settimana e il $19\%$ ogni giorno;
The Conference Board~\cite{conferenceboard2023genai} stima un uso frequente o regolare attorno al $31\%$. Si assume quindi che la frazione di utenti realmente attivi sia compresa
fra il $19\%$ (uso quotidiano) e il $40\%$ (uso frequente) della popolazione abilitabile, con un volume dell'ordine di $30$--$60$ richieste mensili per utente attivo.
Il modello è \emph{robusto} rispetto alle due letture: come mostra la Tab.~\ref{tab:eval-costi}, entrambe convergono su $\sim0{,}55$~milioni di richieste l'anno
e un costo di piattaforma di $\sim$\texteuro$30$~mila.

\begin{table}[htbp]
  \centering
  \caption[Confronto economico a ordini di grandezza]{Confronto annuo fra licenze SaaS nominali e piattaforma a consumo, 
  per la popolazione di $N\approx4\,000$ dipendenti abilitabili di Coop Alleanza. I valori sono parametrici (vedi \texttt{cost\_model.py}) 
  e la frazione di utenti attivi è ancorata ai dati di adozione di Gallup~\cite{gallup2026aiwork}.}
  \label{tab:eval-costi}
  \small
  \setlength{\tabcolsep}{3pt}
  \begin{tabular}{lcr}
    \toprule
    \textbf{Voce} & \textbf{Parametro} & \textbf{Costo annuo (ordine)} \\
    \midrule
    \multicolumn{3}{l}{\emph{Scenario A --- licenze SaaS nominali}} \\
    Licenze per intera popolazione & $N=4\,000 \times {}$\texteuro$25$/mese & $\sim{}$\texteuro$1{,}2\,$M \\
    \midrule
    \multicolumn{3}{l}{\emph{Scenario B --- piattaforma a consumo (realizzata)}} \\
    Utenti realmente attivi & $19$--$40\%$ di $N$ (Gallup) & --- \\
    Richieste/utente attivo/mese & $\sim30$--$60$ & --- \\
    Costo medio per richiesta & $\sim{}$\texteuro$0{,}01$ (routing Flash + agente) & --- \\
    Inferenza LLM (Vertex AI) & $\approx 0{,}55\,$M richieste/anno & $\sim{}$\texteuro$5$--$6$k \\
    Cloud Run + Cloud SQL + base & infrastruttura managed & $\sim{}$\texteuro$25$k \\
    \textbf{Totale piattaforma} & & $\sim{}$\texteuro$30$k \\
    \bottomrule
  \end{tabular}
\end{table}

Con assunzioni prudenti, i due scenari differiscono di circa un \emph{fattore 40} (\texteuro$1{,}2$ milioni contro $\sim$\texteuro$30$ mila all'anno: oltre un ordine di grandezza).
Il vantaggio non deriva dalla tecnologia in sé ma dal \emph{modello di tariffazione}: il SaaS paga la disponibilità per ogni utente abilitato, mentre i dati di adozione
mostrano che l'uso reale si concentra su una minoranza della popolazione. Pagare una licenza per ciascun abilitato significa pagare anche per il $60$--$80\%$ che,
pur avendone accesso, non userebbe lo strumento con regolarità. Il confronto, va ricordato, riguarda il costo \emph{operativo} ricorrente e ignora lo sforzo di sviluppo
e manutenzione della piattaforma, reale e non trascurabile: non è dunque un TCO completo, la cui stima richiederebbe l'ammortamento del lavoro ingegneristico documentato in questa tesi.

% =====================================================================
\section{Discussione}
\label{sec:eval-discussione}
% =====================================================================

\subsection{Risposte alle domande di ricerca}
\label{subsec:eval-rq}

\paragraph{RQ1 --- Agent Accuracy.} Il routing raggiunge il $99{,}2\%$ di accuracy,
sopra il target dell'$85\%$ in entrambe le run.

\paragraph{RQ2 --- Complexity Accuracy.} La stima di complessità passa dall'$88{,}3\%$ (v1.0) al $93{,}3\%$ (v1.1),
restando sopra il target dell'$80\%$. È la metrica che beneficia di più dell'iterazione del prompt ($+5{,}0$~punti percentuali):
le regole aggiunte nella v1.1 hanno reso più stabile la distinzione low/high.

\paragraph{RQ3 --- Calibrazione.} Il confidence score separa nettamente i casi corretti da quello errato:
la confidence media sui corretti è $0{,}922$ in entrambe le run, mentre l'unico caso errato resta ben sotto la soglia di clarification ($0{,}55$ nella v1.0, $0{,}5$ nella v1.1),
rispettando i target dichiarati a priori (Tab.~\ref{tab:eval-metriche}). Il dato va però letto con cautela: poggiando su \emph{un solo} caso errato, la ``media sugli errati'' è di fatto un valore singolo e non una statistica robusta.

\paragraph{RQ4 --- Clarification.} Il clarification rate è risultato molto basso in entrambe le run: il meccanismo è funzionalmente corretto (Fig.~\ref{fig:orion_clarification})
ma scatta troppo di rado per una valutazione statisticamente robusta: servirebbero più casi genuinamente ambigui (confidence 0,50--0,70) e un confidence score più discriminante.

\paragraph{RQ5 --- Iterazione del prompt.} Il passaggio da v1.0 a v1.1 mostra che l'iterazione del prompt è efficace \emph{sulla complessità}:
però non ha inciso a sufficienza sulla correlazione cross-dataset.

\paragraph{RQ6 --- Costo--qualità.} Parzialmente verificata. Sul lato del costo,
la distribuzione della complessità ($86{,}7\%$ \texttt{low}) implica che la grande
maggioranza del traffico è servita dal modello economico, con un risparmio
strutturale che cresce con il rapporto di prezzo fra i modelli fino al limite di
$7{,}5\times$ (Sez.~\ref{subsec:eval-cq-costo}). Sul lato della qualità,
l'execution accuracy non degrada sulla fascia complessa servita dal modello
premium ($9/9$ per il code path Claude Sonnet di SysAid), mentre gli errori si
concentrano sulla fascia economica (Sez.~\ref{subsec:eval-cq-breakdown}).
La quantificazione monetaria del risparmio richiede i volumi di token medi per
richiesta e i listini correnti, non ancora misurati (cfr.\ limiti).

\subsection{Limiti dell'impianto di valutazione}
\label{subsec:eval-limiti}

È importante discutere i limiti dell'impianto sperimentale:

\begin{itemize}
  \item \textbf{Dimensione del gold standard.} 120~test case è un campione sufficiente per identificare
  pattern sistematici di errore ma non per produrre intervalli di confidenza stretti sulle metriche.

  \item \textbf{Bias del costruttore.} I test case sono stati costruiti dall'autore della tesi,
  che pur tentando di essere oggettivo, conosce le capacità e i limiti del sistema.
  Un gold standard costruito da utenti reali senza conoscenza del sistema interno produrrebbe prompt con distribuzione potenzialmente diversa.

  \item \textbf{Dataset sintetico vs traffico reale.} I prompt del gold standard sono sintetici o derivati da log anonimizzati.
  Il traffico reale potrebbe contenere formulazioni, errori ortografici, cambi di lingua o richieste ibride non rappresentati nel gold standard.

  \item \textbf{Non-determinismo intra-run.} Il router è configurato con \texttt{temperature = 0.1}: un valore quasi-deterministico ma non nullo.
  Esecuzioni ripetute dello stesso test case possono produrre classificazioni diverse, specialmente sui casi vicini al confine decisionale.
  Una misura della varianza intra-run ad esempio eseguendo ogni test case $N$~volte e calcolando la percentuale di concordanza quantificherebbe la stabilità del router e
  distinguerebbe gli errori sistematici da quelli stocastici.

  \item \textbf{Quantificazione monetaria del costo--qualità.} Il modello della
  Sez.~\ref{subsec:eval-cq-costo} è parametrico nel rapporto di prezzo $\rho$:
  la conversione in valori monetari richiede i token medi per richiesta
  (misurabili dai \texttt{chat\_logs}) e i listini correnti dei modelli, e
  completerebbe la risposta a RQ6. Analogamente, la scelta di Gemini Flash per
  il router (Sez.~\ref{subsec:orion-model-routing}) è motivata da latenza e
  costo ma non è ancora stata confrontata sperimentalmente con un router su
  modello premium a parità di test suite.

 \end{itemize}
